\documentclass[../main]{subfiles}
\begin{document}
\chapter{Planificación de Proyectos}
\img{images/gantt.jpg}{El diagrama de Gantt es una herramienta
  gráfica cuyo objetivo es exponer el tiempo de dedicación previsto
para diferentes tareas o actividades a lo largo de un tiempo total determinado.}
\info{Contenido}{
  Diagrama de Gantt, Diagrama de PERT, Camino Crítico (CPM).
}
\section{Diagrama de Gantt}
Un diagrama de Gantt es una herramienta esencial para la
planificación de proyectos. Proporciona una vista general de las
tareas programadas, permitiendo que todas las partes implicadas
conozcan el cronograma de ejecución.

Un diagrama de Gantt visualiza:
\begin{enumerate}
  \item La fecha de inicio y finalización del proyecto.
  \item Las tareas desglosadas dentro del proyecto.
  \item La asignación de responsables por tarea.
  \item La duración estimada de cada actividad.
  \item La superposición y relación (dependencias) entre tareas.
\end{enumerate}

\section{Componentes de un diagrama de Gantt}
Normalmente, el diagrama contiene los siguientes elementos:
\begin{enumerate}
  \item \textbf{Eje de tiempo:} Las fechas de inicio y finalización
    permiten visualizar la duración total.
  \item \textbf{Tareas:} Desglose del proyecto en subtareas
    manejables para evitar retrasos u olvidos.
  \item \textbf{Plazos previstos:} Barras horizontales que indican la
    duración temporal de cada tarea.
  \item \textbf{Interdependencias:} Conexiones que indican qué tareas
    deben finalizarse antes de que otras comiencen.
  \item \textbf{Progreso:} Representación visual del porcentaje
    completado de cada tarea respecto a la fecha actual.
\end{enumerate}

\section{Introducción al Método PERT y CPM}
Mientras que el diagrama de Gantt se centra en el \textit{tiempo}, el
diagrama de PERT (\textit{Program Evaluation and Review Technique})
se centra en las \textit{relaciones} entre tareas.

\begin{itemize}
  \item \textbf{Diagrama de Red:} Representación gráfica mediante
    nodos (eventos) y flechas (actividades) que muestra la secuencia lógica.
  \item \textbf{Camino Crítico (CPM):} Es la secuencia de actividades
    que determina la duración más larga del proyecto. Si una
    actividad del camino crítico se retrasa, se retrasa todo el proyecto.
  \item \textbf{Tiempos \textit{Early} y \textit{Last}:}
    \begin{itemize}
      \item \textit{Early} (Temprano): El momento más pronto en que
        una actividad puede comenzar.
      \item \textit{Last} (Tardóo): El momento más tarde en que
        una actividad puede terminar sin retrasar el proyecto.
    \end{itemize}
\end{itemize}

\newpage

\actividad{Guía de Trabajos Prácticos}

\begin{enumerate}
  \item Defina qué es un diagrama de Gantt y mencione tres de sus
    componentes principales.
  \item Explique la diferencia conceptual entre un diagrama de Gantt
    y un diagrama PERT. ¿Cuándo es conveniente utilizar cada uno?

  \item \textbf{Análisis de Proyecto Industrial.}
    Un proyecto presenta las siguientes actividades y dependencias:

    \begin{table}[h!]
      \centering
      \begin{tabular}{|c|c|c|} \hline
        \textbf{Actividad} & \textbf{Duración} & \textbf{Precedente}
        \\ \hline \hline
        A & 2 & - \\
        B & 3 & - \\
        C & 5 & A \\
        D & 2 & B \\
        E & 1 & C, D \\
        F & 3 & E \\
        G & 2 & F \\
        H & 1 & F \\
        I & 4 & G \\
        J & 2 & H, I \\ \hline
      \end{tabular}
      \caption{Tabla de precedencia de tareas - Ejercicio 3}
      \label{tab:ejercicio_3}
    \end{table}

    \begin{enumerate}
      \item Realice el gráfico PERT del proyecto.
      \item Calcule los tiempos \textit{Early} y \textit{Last} en cada nodo.
      \item Identifique y resalte el \textbf{Camino Crítico}.
      \item Represente la gráfica de Gantt correspondiente.
    \end{enumerate}

  \item \textbf{Caso: Clima Cuyo S.R.L.}
    La empresa local se dedica a la climatización de edificios
    mediante sistemas de bomba de calor. Cada instalación requiere
    conectar unidades exteriores e interiores a la red eléctrica. Las
    actividades se detallan a continuación:

    \begin{table}[h!]
      \centering
      \begin{tabular}{|c|l|c|} \hline
        \textbf{Act.} & \textbf{Descripción} & \textbf{Duración (hs)}
        \\ \hline \hline
        A & Estudio, potencia y localización & 4 \\
        B & Instalación de tubos de PVC & 8 \\
        C & Instalación de cables eléctricos & 5 \\
        D & Instalación de interruptores/enchufes & 4 \\
        E & Instalación de consolas interiores & 9 \\
        F & Obras soportes motores exteriores & 7 \\
        G & Instalación de motores exteriores & 5 \\
        H & Conexión de motores y consolas & 0,5 \\
        I & Pruebas y verificación del sistema & 1,5 \\ \hline
      \end{tabular}
      \caption{Actividades Clima Cuyo S.R.L.}
      \label{tab:clima_cuyo}
    \end{table}

    \textbf{Restricciones:}
    \begin{itemize}
      \item La actividad A es predecesora de todas las demás (inicio
        del proyecto).
      \item Las actividades B, C y D son correlativas (secuenciales:
        $B \rightarrow C \rightarrow D$).
      \item Las actividades F y G son correlativas ($F \rightarrow G$).
      \item La actividad E puede realizarse en paralelo a las ramas
        B-C-D y F-G, pero debe comenzar después de A.
      \item La actividad H requiere haber finalizado D, G y E.
      \item La actividad I es la óltima del proceso, realizándose después de H.
    \end{itemize}

    \begin{enumerate}
      \item Elabore el gráfico de Gantt.
      \item Elabore el gráfico PERT.
      \item Indique las actividades del camino crítico y el tiempo de
        holgura de cada nodo.
    \end{enumerate}

  \item \textbf{Caso: Industrias SKIP - Lanzamiento Suave.}
    Para lanzar un nuevo suavizante, se debe instalar nueva
    maquinaria siguiendo estas especificaciones:

    \begin{table}[h!]
      \centering
      \begin{tabular}{|c|l|c|} \hline
        \textbf{Act.} & \textbf{Descripción} & \textbf{Duración (hs)}
        \\ \hline \hline
        A & Derribo de tabiques & 12 \\
        B & Construcción de nuevas divisorias & 20 \\
        C & Nueva instalación eléctrica máquina 1 & 10 \\
        D & Instalación máquina 1 & 24 \\
        E & Nueva instalación eléctrica máquina 2 & 15 \\
        F & Fijación soportes especiales máquina 2 & 6 \\
        G & Instalación máquina 2 & 36 \\
        H & Pintura de los nuevos tabiques & 8 \\
        I & Actividades de prueba del sistema & 5 \\ \hline
      \end{tabular}
      \caption{Actividades Industrias SKIP}
      \label{tab:industrias_skip}
    \end{table}

    \textbf{Secuencia de trabajo:}
    \begin{itemize}
      \item A es la primera actividad, seguida inmediatamente por B.
      \item Al finalizar B, se abren dos ramas paralelas:
        \begin{itemize}
          \item Rama 1: C y D (correlativas: $C \rightarrow D$).
          \item Rama 2: E, F y G (correlativas: $E \rightarrow F
            \rightarrow G$).
        \end{itemize}
      \item Cuando ambas ramas (D y G) finalicen, se procederá con H.
      \item Finalmente, se realizará I.
    \end{itemize}

    \textbf{Consigna:} Realice el gráfico PERT, calcule el tiempo
    total del proyecto e indique el camino crítico.
\end{enumerate}
\end{document}